\begin{proposition}[Invertible covariance with fixed-dimensional marks]
\label{S:marked-theorem}
Let $m_d/d\to c\in(0,\infty)$. Let $(T_i,V_i)_{i\le m_d}$ be i.i.d. with
\begin{equation}
T_i\in[\tau,M]\quad\text{a.s.},\qquad
V_i\in\R^q,\qquad
\E\|V_i\|^r<\infty\quad(r<\infty).
\label{S:marked-assumptions}
\end{equation}
Let $W_i\sim N(0,I_{d-q})$ be i.i.d.\ and independent of the marks, and define
\begin{equation}
Y_i=LV_i+\lambda^{-1/2}L^\perp W_i,\qquad
\widehat M_d=\frac1{m_d}\sum_{i=1}^{m_d}T_iY_iY_i^{\mathsf T}.
\end{equation}
Let $s_c$ be the physical solution of
\begin{equation}
1+zs_c(z)
=
c\,\E\!\left[
\frac{s_c(z)T}{\lambda c+s_c(z)T}
\right],
\label{S:marked-fixed}
\end{equation}
let $\nu_c$ be its measure, and set
\begin{equation}
F_c(z)
=
\lambda c\,\E\!\left[
\frac{T}{\lambda c+s_c(z)T}VV^{\mathsf T}
\right],
\qquad
B_c(z)=zI_q-F_c(z).
\label{S:marked-F}
\end{equation}
Then the empirical spectral distribution of $\widehat M_d$ converges to
$\nu_c$. More precisely, if $K$ is a compact subset of
$(\sup\supp\nu_c,\infty)$, there is a complex neighborhood $U$ of $K$ such
that, with probability tending to one, the orthogonal block $C_{m_d}$ from
\eqref{S:block-matrix} has no spectrum in $U$, $Q_{m_d}$ is holomorphic on
$U$, and on every compact $U'\Subset U$,
\begin{equation}
\sup_{z\in U'}
\left(
\|Q_{m_d}(z)-B_c(z)\|
+
\|Q_{m_d}'(z)-B_c'(z)\|
\right)
\xrightarrow{\Pp}0.
\label{S:Q-convergence}
\end{equation}
Consequently:
\begin{enumerate}
\item eigenvalues of $\widehat M_d$ in compact right intervals are counted by
zeros of $\det B_c$, with multiplicity;
\item root-free compact right intervals contain no eigenvalues with
probability tending to one;
\item on every contour $\Gamma\Subset U$ avoiding the zeros of $\det B_c$,
\begin{equation}
\sup_{z\in\Gamma}
\left\|
L^{\mathsf T}(\widehat M_d-zI)^{-1}L+B_c(z)^{-1}
\right\|
\xrightarrow{\Pp}0;
\label{S:marked-resolvent}
\end{equation}
\item for every eigenpair $(z_d,u_d)$ in such a right region, with
$r_d=L^{\mathsf T}u_d$,
\begin{equation}
\|B_c(z_d)r_d\|=o_{\Pp}(1),
\qquad
\left|r_d^{\mathsf T}B_c'(z_d)r_d-1\right|=o_{\Pp}(1).
\label{S:marked-evec}
\end{equation}
\end{enumerate}
All assertions hold along every deterministic sequence $m_d/d\to c$.
\end{proposition}

\begin{proof}
\emph{Step 1: exact source scaling and independence.}
Let
\begin{equation}
p_d=d-q,
\qquad
\gamma_d=\frac{p_d}{d},
\qquad
c_d=\frac{m_d}{d},
\qquad
\beta_d=\frac{m_d}{p_d}.
\label{S:three-ratios}
\end{equation}
Because $q$ is fixed, $\gamma_d\to1$ and both $c_d$ and $\beta_d$
converge to $c$.  Define
\begin{equation}
\widetilde H_d
=
\frac{\lambda}{d}\sum_{i=1}^{m_d}T_iY_iY_i^{\mathsf T}
=
\lambda c_d\widehat M_d.
\label{S:source-scaling}
\end{equation}
Let $\overline W_d=W/\sqrt{p_d}$ and
$\widetilde W_d=W/\sqrt d=\sqrt{\gamma_d}\,\overline W_d$.  Let
$R_d\in\R^{m_d\times q}$ have $i$th row
$\sqrt{\lambda/d}\,V_i^{\mathsf T}$.  In the basis $[L\ L^\perp]$,
\begin{equation}
\sqrt{\frac{\lambda}{d}}[Y_1\cdots Y_{m_d}]
=
\begin{bmatrix}R_d^{\mathsf T}\\ \widetilde W_d\end{bmatrix}.
\label{S:mark-decomposition}
\end{equation}
Thus $(R_d,D)$ is independent of the Gaussian block
$\overline W_d$.  The rows of $R_d$ may be correlated with $D$ because
both are functions of the joint mark $(T,V)$; this causes no problem because
the source anisotropic theorem permits random test vectors independent of
the Gaussian block.  Gaussianity of $V$ is never used.

The distinction between $d$ and $p_d=d-q$ is kept exact.  Put
\begin{equation}
D_d^\circ=\gamma_dD.
\label{S:rescaled-D}
\end{equation}
Then the orthogonal block of $\widetilde H_d$ is exactly
\begin{equation}
\widetilde C_d
=
\widetilde W_dD\widetilde W_d^{\mathsf T}
=
\overline W_dD_d^\circ\overline W_d^{\mathsf T}.
\label{S:orthogonal-exact-scaling}
\end{equation}
Accordingly, the source local law is applied with row dimension $p_d$,
aspect ratio $\beta_d=m_d/p_d$, and diagonal weights $D_d^\circ$.
No finite-dimensional identification of $d$ with $d-q$ is made.

\emph{Step 2: verify the imported weighted-Wishart hypotheses.}
Let $\rho$ be the law of $T$.  The empirical law
$\widehat\rho_d=m_d^{-1}\sum_i\delta_{T_i}$ converges weakly to $\rho$
almost surely.  Its support converges to $\supp\rho$ in Hausdorff distance:
for every $\eta>0$, cover the compact support by finitely many open
$(\eta/2)$-balls centered in the support.  Each ball has positive $\rho$-mass,
so the probability that it is missed is exponentially small in $m_d$; a
union bound and Borel--Cantelli apply.  The reverse inclusion is automatic
because $T_i\in\supp\rho$ almost surely.

The empirical law of $D_d^\circ$ is the dilation of
$\widehat\rho_d$ by $\gamma_d$. If one insists on the source paper's
literal low-dimensional representation of a diagonal entry, any law on
$[\tau,M]$ can be realized as
$T=F_T^{-1}(\Phi(\xi))$ for a standard Gaussian $\xi$ and a generalized
quantile $F_T^{-1}$. Thus the formal coupling assumption for the independent
weighted-Wishart input entails no additional restriction here. Since
$\gamma_d\to1$, the empirical law converges weakly
to $\rho$, and its support converges to $\supp\rho$ in Hausdorff distance.
For all sufficiently large $d$ its entries belong to
$[\tau/2,2M]$.  These are the strong empirical convergence and regularity
hypotheses used in Proposition~2.8 and Theorem~2.12 of
Ref.~\cite{BenArousEtAl2025}.  For a prescribed bounded spectral compact, choose the source regularity
parameter smaller than
\begin{equation}
\min\{\tau/2,(2M)^{-1},(1+\sup|z|)^{-1},1\}.
\end{equation}
The weight law is then regular with this parameter, and the compact lies
inside the spectral domain of the source estimate.
The source no-outlier theorem therefore applies to \eqref{S:orthogonal-exact-scaling}.

The source anisotropic estimate also applies to the finitely many columns
of $R_d$.  They are independent of $\overline W_d$, may be correlated with
$D$, and for every fixed $j\le q$,
\begin{equation}
\|R_d e_j\|^2
=
\frac{\lambda}{d}\sum_{i=1}^{m_d}V_{ij}^2
=O(1)
\label{S:R-column-bound}
\end{equation}
almost surely on the natural coupling, by the strong law and
$m_d/d\to c$.  The random-vector clause in
Proposition~\ref{S:source-input} is used directly: each nonzero vector
$R_d e_j$ is normalized and embedded in the lower block of the source
linearization.  To recover every matrix entry, apply the same estimate to
the finite collection $R_d e_i$, $R_d e_j$, and
$R_d(e_i\pm e_j)$ and use polarization; zero vectors contribute
identically zero.  Because $q$ is fixed, a finite union of the corresponding
high-probability events still has probability tending to one.  Restoring
the $O(1)$ column norms in \eqref{S:R-column-bound} preserves convergence.
Thus no quenched assertion beyond the source theorem is being used, and no
Gaussianity assumption on the joint mark $(T,V)$ is introduced.

\emph{Step 3: the exact Schur matrix and the local law.}
For $z\notin\sigma(\widetilde C_d)$, the source-scaled Schur matrix is
\begin{equation}
\widetilde Q_d(z)
=
zI_q-R_d^{\mathsf T}DR_d
+
R_d^{\mathsf T}D\widetilde W_d^{\mathsf T}
(\widetilde C_d-zI)^{-1}
\widetilde W_dDR_d.
\label{S:source-Q}
\end{equation}
Equivalently,
\begin{equation}
\widetilde Q_d(z)=zI_q-R_d^{\mathsf T}\mathcal K_d(z)R_d,
\quad
\mathcal K_d(z)
=
D-D\widetilde W_d^{\mathsf T}
(\widetilde C_d-zI)^{-1}\widetilde W_dD.
\label{S:K-block}
\end{equation}
Let
\begin{equation}
\mathcal K_d^\circ(z)
=
D_d^\circ-D_d^\circ\overline W_d^{\mathsf T}
(\widetilde C_d-zI)^{-1}\overline W_dD_d^\circ.
\end{equation}
Using $D_d^\circ=\gamma_dD$ and
$\widetilde W_d=\sqrt{\gamma_d}\,\overline W_d$ gives the exact identity
\begin{equation}
\mathcal K_d^\circ(z)=\gamma_d\mathcal K_d(z).
\label{S:K-gamma-identity}
\end{equation}

Let $\overline s_d$ be the source deterministic transform associated with
the empirical law of $D_d^\circ$ and aspect ratio $\beta_d$.  The
lower-right block of the source anisotropic local law gives, locally
uniformly on every compact set separated from the limiting support,
\begin{equation}
\mathcal K_d^\circ(z)
\ \sim_{\rm aniso}\
D_d^\circ(I+\overline s_d(z)D_d^\circ)^{-1}.
\end{equation}
Set
\begin{equation}
a_d(z)=\gamma_d\overline s_d(z).
\label{S:a-d-definition}
\end{equation}
After division by $\gamma_d$ in \eqref{S:K-gamma-identity},
\begin{equation}
\mathcal K_d(z)
\ \sim_{\rm aniso}\
D(I+a_d(z)D)^{-1}.
\label{S:K-local-law}
\end{equation}
The finite-dimensional deterministic equation is worth recording.  Since
$D_d^\circ=\gamma_dD$, equation
\eqref{S:source-input-fixed} for $\overline s_d$ becomes, after multiplying
by $\gamma_d^{-1}$,
\begin{equation}
\frac1{a_d(z)}
=-\frac{z}{\gamma_d}
+\beta_d\frac1{m_d}\sum_{i=1}^{m_d}
\frac{T_i}{1+a_d(z)T_i}.
\label{S:a-d-finite-equation}
\end{equation}
The empirical laws of $D_d^\circ$ converge strongly to $\rho$,
$\gamma_d\to1$, and $\beta_d\to c$.  Lemma~2.13 of the source, or directly
the stability of \eqref{S:a-d-finite-equation} on compact sets separated
from the limiting support, therefore gives
\begin{equation}
a_d\longrightarrow\widetilde s
\label{S:a-d-limit}
\end{equation}
locally uniformly, where $\widetilde s$ is the physical solution in
\eqref{S:source-fixed} below.

Applying \eqref{S:K-local-law} to the normalized columns of $R_d$, using
\eqref{S:R-column-bound}, and polarizing the resulting quadratic-form
estimates yields
\begin{equation}
\widetilde Q_d(z)
=
zI_q-R_d^{\mathsf T}D(I+a_d(z)D)^{-1}R_d+o_{\Pp}(1)
\label{S:Q-local-law}
\end{equation}
locally uniformly.  This is the precise point at which the source local law
is used.

\emph{Step 4: denominator separation and the marked law of large numbers.}
The limiting transform in the source scaling solves
\begin{equation}
1+z\widetilde s(z)
=
c\,\E\!\left[
\frac{\widetilde s(z)T}{1+\widetilde s(z)T}
\right].
\label{S:source-fixed}
\end{equation}
Let $M_\rho=\sup\supp\rho$.  On the real component to the right of the
limiting support, the Silverstein--Choi support parametrization
\cite{SilversteinChoi1995}, also used in the proof of Proposition~2.8 of
Ref.~\cite{BenArousEtAl2025}, states that
\begin{equation}
y(z):=-\frac1{\widetilde s(z)}>M_\rho,
\qquad
z=y+c\,y\int\frac{t}{y-t}\,\rho(dt).
\label{S:support-map}
\end{equation}
Consequently, for $t\in\supp\rho$,
\begin{equation}
1+\widetilde s(z)t=1-\frac{t}{y(z)}>0.
\end{equation}
If $K$ is a compact right set, continuity of $y$ and compactness give a
number $\delta_K>0$ such that
\begin{equation}
\inf_{\substack{z\in K\\t\in\supp\rho}}
|1+\widetilde s(z)t|\ge\delta_K.
\label{S:marked-denominator}
\end{equation}
Holomorphy and uniform continuity preserve half this bound on a sufficiently
thin complex neighborhood $U$ of $K$.  Strong support convergence and
\eqref{S:a-d-limit} then give, with probability tending to one,
\begin{equation}
\inf_{\substack{z\in U'\\1\le i\le m_d}}
|1+a_d(z)T_i|\ge\delta_K/4
\label{S:finite-denominator}
\end{equation}
for every compact $U'\Subset U$ and all sufficiently large $d$.

The reduced matrix in \eqref{S:Q-local-law} is
\begin{equation}
\widetilde F_d(z)
=
R_d^{\mathsf T}D(I+a_d(z)D)^{-1}R_d
=
\frac{\lambda}{d}\sum_{i=1}^{m_d}
\frac{T_iV_iV_i^{\mathsf T}}{1+a_d(z)T_i}.
\label{S:empirical-F}
\end{equation}
Insert the deterministic transform $\widetilde s$ to separate the three
sources of error:
\begin{align}
\widetilde F_d(z)
-\lambda c\,\E\!\left[
\frac{T}{1+\widetilde s(z)T}VV^{\mathsf T}
\right]
={}&\frac{\lambda}{d}\sum_{i=1}^{m_d}T_iV_iV_i^{\mathsf T}
\left[
\frac1{1+a_d(z)T_i}
-\frac1{1+\widetilde s(z)T_i}
\right]\notag\\
&+\lambda c_d\left\{
\frac1{m_d}\sum_{i=1}^{m_d}
\frac{T_iV_iV_i^{\mathsf T}}{1+\widetilde s(z)T_i}
-\E\!\left[
\frac{T}{1+\widetilde s(z)T}VV^{\mathsf T}
\right]
\right\}\notag\\
&+\lambda(c_d-c)\,
\E\!\left[
\frac{T}{1+\widetilde s(z)T}VV^{\mathsf T}
\right].
\label{S:marked-LLN-decomposition}
\end{align}
The first term is uniformly $o(1)$ on $U'$ by
\eqref{S:a-d-limit}, \eqref{S:finite-denominator}, and
$m_d^{-1}\sum_i\|V_i\|^2=O(1)$ almost surely.  For the second term, every
matrix entry has an integrable envelope $C(1+\|V\|^2)$.  A strong law on a
finite net in $U'$, followed by the derivative bound implied by
\eqref{S:marked-denominator}, gives a uniform strong law.  The final term
vanishes because $c_d\to c$.  Therefore
\begin{equation}
\sup_{z\in U'}
\left\|
\widetilde F_d(z)
-
\lambda c\,\E\!\left[
\frac{T}{1+\widetilde s(z)T}VV^{\mathsf T}
\right]
\right\|
\longrightarrow0
\quad\text{almost surely}.
\label{S:uniform-mark-LLN}
\end{equation}
Combining \eqref{S:Q-local-law} and \eqref{S:uniform-mark-LLN} gives locally
uniform convergence of $\widetilde Q_d$ to
\begin{equation}
\widetilde B(z)
=
zI_q-
\lambda c\,\E\!\left[
\frac{T}{1+\widetilde s(z)T}VV^{\mathsf T}
\right].
\label{S:B-tilde}
\end{equation}
All functions are holomorphic on the same neighborhood.  Cauchy's formula
therefore gives locally uniform convergence of their first derivatives as
well.

\emph{Step 5: return to the $1/m_d$ normalization.}
The limiting deterministic rescaling is
\begin{equation}
\widetilde s(\lambda cx)=\frac1{\lambda c}s_c(x),
\qquad
\widetilde B(\lambda cx)=\lambda c\,B_c(x).
\label{S:marked-rescaling}
\end{equation}
At finite dimension, the exact identity
$\widetilde H_d=\lambda c_d\widehat M_d$ and the block definitions give
\begin{equation}
\widetilde Q_d(\lambda c_dx)=\lambda c_d\,Q_{m_d}(x).
\label{S:finite-Q-scaling}
\end{equation}
This is the only finite-dimensional scaling used. The source transform
contains the separate ratio $\beta_d=m_d/(d-q)$, and
$\beta_d-c_d\to0$ because $q$ is fixed. Since both ratios converge to $c$,
the locally uniform convergence just proved implies
\eqref{S:Q-convergence}.

Finally, $\widetilde H_d$ differs from
$\diag(0_q,\widetilde C_d)$ by rank at most $2q$, so their empirical
spectral measures, normalized by $d$, have the same limit.  The difference
between normalization by $d$ and by $p_d=d-q$ is $o(1)$.  Since
$\widetilde H_d=\lambda c_d\widehat M_d$ and $c_d\to c$, deterministic
dilation of that limiting measure gives the measure $\nu_c$ whose transform
satisfies \eqref{S:marked-fixed}.  This proves the bulk claim in the original
$1/m_d$ normalization.

We now justify the zero count without hiding any complex roots.  For
$z\in\C_+$, the physical transform satisfies
$\operatorname{Im}\widetilde s(z)>0$, and for every $T>0$,
\begin{equation}
\operatorname{Im}\frac{T}{1+\widetilde s(z)T}
=
-\frac{T^2\operatorname{Im}\widetilde s(z)}
{|1+\widetilde s(z)T|^2}<0.
\end{equation}
Hence $\operatorname{Im}\widetilde B(z)\succeq
(\operatorname{Im}z)I_q\succ0$.  By conjugation the analogous statement
holds in $\C_-$, so every zero of $\det\widetilde B$ off the limiting bulk
is real.  The same conclusion holds for $B_c$ by
\eqref{S:marked-rescaling}.

Let $[a,b]$ be a compact right interval whose endpoints are not limiting
roots.  Choose a thin rectangle $\mathcal R$ with vertical sides through
$a$ and $b$, contained in the common holomorphy domain.  Its horizontal
sides contain no zero by the preceding imaginary-part argument, and its
vertical sides contain no zero after the height is chosen sufficiently
small.  Therefore
\begin{equation}
\inf_{z\in\partial\mathcal R}|\det B_c(z)|>0.
\end{equation}
Uniform convergence in \eqref{S:Q-convergence} implies uniform convergence
of the determinants on $\partial\mathcal R$; Rouch\'e's theorem, equivalently
the argument principle, gives the same number of zeros inside
$\mathcal R$, counted analytically.  The source no-outlier theorem gives no
orthogonal-block eigenvalue in $[a,b]$ with probability tending to one.
Since the spectrum of $C_{m_d}$ is real, its determinant has no zero inside
$\mathcal R$.  The exact factorization \eqref{S:exact-det} then identifies
the zeros of $\det Q_{m_d}$ inside $\mathcal R$ with the eigenvalues of the
real symmetric full matrix in $[a,b]$, counted algebraically.  This proves
the count and root-free assertions, including multiplicity.

On a contour avoiding limiting roots, uniform convergence of $Q_{m_d}$ and
the matrix inverse identity \eqref{S:exact-compressed-resolvent} give
\eqref{S:marked-resolvent}.  Finally,
\eqref{S:exact-evec-Q} and \eqref{S:exact-evec-norm}, together with uniform
convergence of $Q_{m_d}$ and $Q_{m_d}'$, give
\eqref{S:marked-evec}.  This proves every assertion of the proposition.
\end{proof}

